\section{理论与假设}

\begin{frame}{理论基础：三条逻辑共同支撑研究框架}
  \begin{columns}[T,onlytextwidth]
    \column{0.32\textwidth}
    \begin{block}{外部性理论}
      碳排放具有典型负外部性，市场会对高碳活动 \hl{定价不足}。
    \end{block}

    \column{0.32\textwidth}
    \begin{exampleblock}{金融功能理论}
      金融体系通过 \hlb{资源配置、风险管理、信息筛选} 改变资金流向。
    \end{exampleblock}

    \column{0.32\textwidth}
    \begin{alertblock}{波特假说}
      合理环境约束与金融激励可推动企业技术升级与绿色创新。
    \end{alertblock}
  \end{columns}

  \vspace{0.8em}
  \begin{itemize}
    \item 理论上，绿色金融应通过融资约束与资源倾斜，降低高碳活动比重。
    \item 但不同地区资源禀赋不同，实际效果可能存在显著 \hl{区域异质性}。
  \end{itemize}
\end{frame}

\begin{frame}{研究假设与分析思路}
  \begin{block}{假设 H1}
    绿色金融发展水平提升，将显著降低各地区碳排放强度。
  \end{block}

  \begin{exampleblock}{假设 H2}
    绿色金融的减排效应具有区域异质性，在不同地区的边际作用不同。
  \end{exampleblock}

  \vspace{0.6em}
  \begin{itemize}
    \item 如果 H1 成立，说明绿色金融不是概念性政策，而是有可测量的减排效果。
    \item 如果 H2 成立，说明政策设计不能“一刀切”，而应强调 \hlb{区域适配性}。
  \end{itemize}
\end{frame}
